\section{Gabarito e resoluções comentadas — Língua Portuguesa}

\begin{solucao}{01.01}\textbf{CERTO.} \emph{Embora} marca concessão; \emph{porque}, causa/explicação.\end{solucao}
\begin{solucao}{01.02}\textbf{ERRADO.} ``Nem todos'' nega a universalidade; ``nenhum'' afirma quantidade zero.\end{solucao}
\begin{solucao}{01.03}\textbf{CERTO.} \emph{Seu} pode retomar gestor ou auditor.\end{solucao}
\begin{solucao}{01.04}\textbf{CERTO.} \emph{Haver} existencial é impessoal; \emph{existir} concorda: ``existem inconsistências''.\end{solucao}
\begin{solucao}{01.05}\textbf{ERRADO.} O correto é ``deve haver soluções''.\end{solucao}
\begin{solucao}{01.06}\textbf{CERTO.} A oração passa de restritiva a explicativa, alterando o alcance.\end{solucao}
\begin{solucao}{01.07}\textbf{CERTO.} \emph{Referir-se a} + artigo \emph{a} produz crase.\end{solucao}
\begin{solucao}{01.08}\textbf{CERTO.} A negação atrai o pronome: ``não se verificaram''.\end{solucao}
\begin{solucao}{01.09}\textbf{ERRADO.} \emph{Retificar}=corrigir; \emph{ratificar}=confirmar.\end{solucao}
\begin{solucao}{01.10}\textbf{CERTO.} ``Se viesse... concluiríamos'' apresenta correlação verbal adequada.\end{solucao}
\begin{solucao}{01.11}\textbf{B.} O trecho autoriza apenas concluir que uma métrica melhorou e outra piorou.\end{solucao}
\begin{solucao}{01.12}\textbf{B.} \emph{Pois} explica a razão da devolução.\end{solucao}
\begin{solucao}{01.13}\textbf{CERTO.} \emph{Caso}=condição; \emph{embora}=concessão.\end{solucao}
\begin{solucao}{01.14}\textbf{C.} O adjunto deslocado está corretamente separado por vírgula.\end{solucao}
\begin{solucao}{01.15}\textbf{B.} \emph{Existir} é pessoal: ``devem existir alternativas''.\end{solucao}
\begin{solucao}{01.16}\textbf{D.} É correta a construção ``informar algo a alguém''.\end{solucao}
\begin{solucao}{01.17}\textbf{CERTO.} Admitem-se ``informar algo a alguém'' e ``informar alguém de algo''.\end{solucao}
\begin{solucao}{01.18}\textbf{D.} \emph{Dirigir-se a} + \emph{a unidade} = \emph{à unidade}.\end{solucao}
\begin{solucao}{01.19}\textbf{B.} A negação exige/prioriza próclise no registro formal.\end{solucao}
\begin{solucao}{01.20}\textbf{CERTO.} \emph{Que} retoma \emph{relatório}; é pronome relativo.\end{solucao}
\begin{solucao}{01.21}\textbf{B.} A oração completa o verbo \emph{afirmou}.\end{solucao}
\begin{solucao}{01.22}\textbf{A.} Preserva agente, sequência e conteúdo básico.\end{solucao}
\begin{solucao}{01.23}\textbf{CERTO.} A voz muda, mas os participantes básicos permanecem.\end{solucao}
\begin{solucao}{01.24}\textbf{C.} \emph{Por que} está correto na interrogativa indireta.\end{solucao}
\begin{solucao}{01.25}\textbf{CERTO.} Tempo decorrido pede \emph{há}.\end{solucao}
\begin{solucao}{01.26}\textbf{B.} \emph{Somente... inicialmente} preserva a restrição temporal de \emph{apenas}.\end{solucao}
\begin{solucao}{01.27}\textbf{CERTO.} ``Nenhum'' é logicamente mais forte que ``ao menos um não''.\end{solucao}
\begin{solucao}{01.28}\textbf{B.} I restringe; II explica e pode atribuir a característica ao conjunto.\end{solucao}
\begin{solucao}{01.29}\textbf{B.} \emph{Apesar de} preserva a concessão.\end{solucao}
\begin{solucao}{01.30}\textbf{CERTO.} \emph{Ainda que} seleciona naturalmente o subjuntivo.\end{solucao}
\begin{solucao}{01.31}\textbf{A.} \emph{Aspirar a} + \emph{a melhoria} = \emph{à melhoria}.\end{solucao}
\begin{solucao}{01.32}\textbf{C.} ``Deste'' torna o analista o referente; em redação real, ``parecer do analista'' seria ainda mais direto.\end{solucao}
\begin{solucao}{01.33}\textbf{CERTO.} \emph{Fazer} temporal é impessoal: ``faz cinco meses''.\end{solucao}
\begin{solucao}{01.34}\textbf{B.} Com artigo definido: ``é necessária a revisão''.\end{solucao}
\begin{solucao}{01.35}\textbf{A.} Há anáfora e relação adversativa.\end{solucao}
\begin{solucao}{01.36}\textbf{CERTO.} \emph{Portanto} introduz conclusão.\end{solucao}
\begin{solucao}{01.37}\textbf{B.} \emph{-la} retoma \emph{a norma}.\end{solucao}
\begin{solucao}{01.38}\textbf{C.} É a única redação formal e objetiva.\end{solucao}
\begin{solucao}{01.39}\textbf{CERTO.} A troca transforma possibilidade em certeza futura.\end{solucao}
\begin{solucao}{01.40}\textbf{B.} O sujeito elíptico da oração reduzida é a equipe.\end{solucao}
\begin{solucao}{01.41}\textbf{ERRADO.} \emph{Prioritariamente} indica preferência/ordem, não exclusividade.\end{solucao}
\begin{solucao}{01.42}\textbf{B.} Preserva concessão e a diferença entre controle formal e operação comprovada.\end{solucao}
\begin{solucao}{01.43}\textbf{B.} Nega-se ``todos'': pelo menos um irregular não foi identificado.\end{solucao}
\begin{solucao}{01.44}\textbf{CERTO.} Possibilidade e certeza têm valores modais distintos.\end{solucao}
\begin{solucao}{01.45}\textbf{D.} \emph{Embora} troca causa por concessão.\end{solucao}
\begin{solucao}{01.46}\textbf{CERTO.} A subordinada com \emph{que} favorece ``que se revisassem''.\end{solucao}
\begin{solucao}{01.47}\textbf{C.} É mensurável, específica e separa fato de opinião.\end{solucao}
\begin{solucao}{01.48}\textbf{B.} A expressão causal intercalada fica isolada.\end{solucao}
\begin{solucao}{01.49}\textbf{CERTO.} As vírgulas marcam o segmento intercalado entre sujeito e predicado.\end{solucao}
\begin{solucao}{01.50}\textbf{C.} Elimina coloquialismo, ambiguidade e imprecisão.\end{solucao}
